\documentclass[10pt,a4paper]{amsart}
\usepackage[margin=19mm]{geometry}
\usepackage{amsmath,amssymb,mathtools}
\usepackage[scr=rsfs,cal=euler]{mathalfa}
\usepackage{xfrac}
\usepackage[unicode,hidelinks,bookmarks=false]{hyperref}
\newcommand{\ie}{i.e.\ }
\newcommand{\cf}{cf.\ }
\newcommand{\resp}{respectively\ }
\newcommand{\Subsection}{\subsection}
\providecommand{\nobreakdash}{\nobreak\hbox{-}}
\setlength{\emergencystretch}{2em}
\multlinegap0pt
\usepackage{bm}
\usepackage{bbm}
\usepackage{dsfont}
\usepackage{xspace}
\usepackage{cancel}
\usepackage{mathtools}
\usepackage{amsthm}
\makeatletter
\@ifundefined{theorem}{\newtheorem{theorem}{Theorem}}{}
\@ifundefined{lemma}{\newtheorem{lemma}[theorem]{Lemma}}{}
\makeatother
\title[Monoidal Categorification and Quantization of Braid Varietie]{Monoidal Categorification and Quantization of Braid Varieties}
\author{Yingjin Bi}
\date{Source: arXiv:2609.02962v2 (CC BY 4.0)}
\begin{document}
\maketitle

\noindent\emph{Source and licence.} Monoidal Categorification and Quantization of Braid Varieties. Version arXiv:2609.02962v2 (CC BY 4.0), distributed under the Creative Commons Attribution 4.0 International licence, as recorded in the arXiv licence metadata for this exact version and shown on the abstract page https://arxiv.org/abs/2609.02962v2. Full rights notice: licenses/arxiv-2609.02962-CC-BY-4.0.txt.\par\smallskip

\noindent\textit{Editorial scope.} Counted unit: the Lemma whose source label is \texttt{lem:gamma-vc-well-defined} in the article (source0.tex lines 6011--6015, source label \texttt{lem:gamma-vc-well-defined}), delivered together with the article's own complete original proof of it (source0.tex lines 6017--6138, one \texttt{proof} environment of 122 lines). No mathematics is written, repaired or re-derived: statement and proof are the article's own text line for line. Required context retained uncounted: none. Every definition, display and label of those lines is retained unchanged. Omitted from this package and not redistributed: 1--6010, 6016--6016, 6139--11445. Nothing inside a retained excerpt is omitted or altered: every retained body is delivered exactly as the article's lines are written, so no omission receipt of any kind accompanies this package. Citations: the retained excerpts carry no citation command at all, so no bibliography entry of the article is transcribed and no reference is redistributed. Reference adaptation: none was needed and none was made; every reference inside the delivered unit resolves inside it, so no unresolved reference or citation is produced. Printed slips kept literal: no printed word, symbol, hypothesis or number of the counted statement or of its proof was altered. Layout and class: the article's own class is \texttt{amsart} with packages including mathrsfs, amsfonts, amsmath, amssymb, array, amsthm, graphicx, caption, enumerate, enumitem, geometry, tikz-cd, tikz, rotating; the wrapper is the lane's pinned \texttt{amsart} editorial wrapper at 10pt on A4, which restates the theorem-like environments the retained text needs and adds the article's own macro definitions that the retained text needs (none), with extra wrapper packages (bm, bbm, dsfont, xspace, cancel, mathtools). The delivered numbering is the unit's own. Assets: no retained span contains a figure environment, a graphics-inclusion command, a PDF-page inclusion, a table or a TikZ picture; no image file is copied into this package, the package declares no asset dependency and no image file is redistributed with it. Licence: version arXiv:2609.02962v2 of this article is distributed under the Creative Commons Attribution 4.0 International licence, as recorded in the arXiv licence metadata for this exact version and shown on the abstract page https://arxiv.org/abs/2609.02962v2. A full rights notice is delivered at licenses/arxiv-2609.02962-CC-BY-4.0.txt. Characters: every retained excerpt is pure ASCII, so the delivered engine, XeLaTeX with its default Latin Modern text font, reports no missing character.
\begin{lemma}
\label{lem:gamma-vc-well-defined}
The coroot \(\gamma_{v,c}\) is independent of the chosen reduced
expression of \(w_c\).
\end{lemma}

\begin{proof}
By Matsumoto's theorem, any two reduced expressions of \(w_c\) are
related by a sequence of braid moves. Since \(G\) is simply laced, it
suffices to consider commutation moves and braid moves of the form
\[
    iji\longleftrightarrow jij.
\]

A commutation move interchanges two orthogonal simple reflections.
The corresponding two inversion coroots and the two cycle
coefficients are interchanged simultaneously, so their contribution
to \(\gamma_{v,c}\) is unchanged.

Consider a \(3\)-move
\[
    \mathbf j=\mathbf x(iji)\mathbf y
    \quad\longleftrightarrow\quad
    \mathbf j'=\mathbf x(jij)\mathbf y,
\]
where \(i\) and \(j\) are adjacent. Write
\(
    \mathbf y=(y_1,\ldots,y_q)
\)
and let
\(
    u:=s_{y_q}\cdots s_{y_1}.
\)
Thus \(u\) is the common Weyl group action on the coroot lattice
coming from the suffix \(\mathbf y\).

Before applying \(u\), the three inversion coroots associated with
the substring \(iji\) are
\[
    \bigl(
        \alpha_j^\vee,\,
        \alpha_i^\vee+\alpha_j^\vee,\,
        \alpha_i^\vee
    \bigr),
\]
whereas those associated with \(jij\) are
\[
    \bigl(
        \alpha_i^\vee,\,
        \alpha_i^\vee+\alpha_j^\vee,\,
        \alpha_j^\vee
    \bigr).
\]
Consequently, the actual inversion coroot triples in the full words
\(\mathbf j\) and \(\mathbf j'\) are respectively
\[
    \bigl(
        u(\alpha_j^\vee),\,
        u(\alpha_i^\vee+\alpha_j^\vee),\,
        u(\alpha_i^\vee)
    \bigr)
\]
and
\[
    \bigl(
        u(\alpha_i^\vee),\,
        u(\alpha_i^\vee+\alpha_j^\vee),\,
        u(\alpha_j^\vee)
    \bigr).
\]

Let \((a_1,a_2,a_3)\) be the values of the Lusztig cycle on the three
affected edges. Under the braid move, they become
\[
\begin{aligned}
    b_1=a_2+a_3-\min(a_1,a_3),\quad
    b_2=\min(a_1,a_3),\quad
    b_3=a_1+a_2-\min(a_1,a_3).
\end{aligned}
\]
In particular,
\[
    b_1+b_2=a_2+a_3,
    \qquad
    b_2+b_3=a_1+a_2.
\]
Hence
\begin{align*}
    &a_1\alpha_j^\vee
     +a_2(\alpha_i^\vee+\alpha_j^\vee)
     +a_3\alpha_i^\vee\\
    &\qquad
     =(a_2+a_3)\alpha_i^\vee
       +(a_1+a_2)\alpha_j^\vee\\
    &\qquad
     =b_1\alpha_i^\vee
       +b_2(\alpha_i^\vee+\alpha_j^\vee)
       +b_3\alpha_j^\vee.
\end{align*}
Applying the common linear action \(u\) to this equality gives
\begin{align*}
    &a_1u(\alpha_j^\vee)
     +a_2u(\alpha_i^\vee+\alpha_j^\vee)
     +a_3u(\alpha_i^\vee)\\
    &\qquad
     =b_1u(\alpha_i^\vee)
       +b_2u(\alpha_i^\vee+\alpha_j^\vee)
       +b_3u(\alpha_j^\vee).
\end{align*}
Thus the contribution of the three affected positions to
\(\gamma_{v,c}\) is unchanged.

The coefficients outside these three positions are unchanged by the
local tropical braid move. The inversion coroots associated with
positions in the suffix \(\mathbf y\) are plainly unchanged. Those
associated with positions in the prefix \(\mathbf x\) are also
unchanged, because their defining suffix actions contain respectively
\[
    s_is_js_i
    \quad\text{and}\quad
    s_js_is_j,
\]
which are equal by the braid relation. Therefore every contribution
outside the affected triple is unchanged as well.

Hence \(\gamma_{v,c}\) is invariant under a \(3\)-move, and therefore
under every change of reduced expression.
\end{proof}

\end{document}
